\cotefile{53}
\scanpage{16}
\[
\left\{\begin{aligned}
&\text{Groups over bases } S\ \text{of any kind. Take base } S\ \deleted{\unsure{?}}\ \text{or take base } k\ \text{?}\\
&\text{Non-affine groups O.K.}
\end{aligned}\right.
\]

Linear representations of the $D(M)$ [special cases: $\mathbb{G}_m$, $\mu_n$]

Application: homomorphisms $D(M)\longrightarrow\mathbb{G}_m$: one finds that $M_S$ is ``reflexive'' with respect to $\mathbb{G}_m$.

\uline{Cartier duality}\quad $\underline{\mathrm{Hom}}_{\mathrm{gr}}(G,\mathbb{G}_m)$ \deleted{\unsure{$\overset{\Delta}{=}$|is}} representable? It is if $G$ diagonalizable, but \deleted{not}\inserted{in general (= for $G=\mathbb{G}_m$)} for an affine group scheme (even of finite type). In contrast, if $G$ is finite locally free over $S$, one has
% slanted
\reversemarginpar\margin{$\underline{\mathrm{Hom}}_{\mathrm{gr}}(G,H)$ representable if $G$ is finite locally free, and $H$ affine (say)}\normalmarginpar
\[
\underline{\mathrm{Hom}}_{\mathrm{gr}}(G,\mathbb{G}_m)
\subset \mathrm{Hom}_{\unsure{\mathcal O|O|\mathbb O}\text{-alg}}
(\mathcal O_S[T,T^{-1}],\mathcal A)\simeq W(\mathcal A)^*,
\]
and the subfunctor is that of the \deleted{$g$} $u\in W(\mathcal A)^*(S')$ such that
\[
\Pi(u)=u\otimes u\qquad \varepsilon(u)=1
\]
\margin{[\,\uline{NB} One can therefore forget the condition $u$ invertible since $\underline{\mathrm{Hom}}_{\mathrm{gr}}(G,\mathbb{G}_m)=\underline{\mathrm{Hom}}_{\mathrm{mon}}(G,\mathcal O^{\times})$.}

\deleted{The} The representability is clear, but moreover one recognizes that it is \deleted{\unsure{the|at}} $\mathrm{Spec}(\check{\mathcal A})$, $\check{\mathcal A}$ \deleted{being}\inserted{$=\underline{\mathrm{Dist}}(G)_{S}$} an algebra thanks to the coalgebra structure of $\mathcal A$. On the other hand, the coalgebra structure of $\check{\mathcal A}$ (coming from the fact that the functor it represents is a group functor
% slanted
\reversemarginpar\margin{$\check{\mathcal A}$ commutative since $G$ is. If $G$ is not commutative, one must take the functor $k'\longmapsto\mathrm{Hom}_{k\text{-alg}}(\check{\mathcal A},k')=\mathrm{Hom}_{k\text{-alg}}({\check{\mathcal A}\deleted{\unsure{)}}}_{\mathrm{ab}},k')$}\normalmarginpar

\uline{Theorem} Every finite locally free commutative group \inserted{$G$} is reflexive for
\par\noindent\begin{tabular}{@{}|p{.91\linewidth}@{}}
$\mathbb{G}_m$. Its dual $G^*$ is finite locally free commutative, hence an \uline{\unsure{of involution|involution}} $G\unsure{\mathrel{\vert\!-\!{\cdot}}|\mapsto}G^*$ of the category of locally free commutative groups. It corresponds to the duality of bicommutative bialgebras over $S$.
\end{tabular}

\[
\begin{aligned}
\text{\uline{Examples}}\qquad &\mu_n\longleftrightarrow\mathbb Z/n\mathbb Z
&\ \Bigg|\, &\text{But also }\ D(M)\longleftrightarrow M_S\\
&\alpha_p\longleftrightarrow\alpha_p && 
\end{aligned}
\]

Works $\pm$ for \deleted{affine} group schemes \deleted{\unsure{to|has}} over $S$ \deleted{such that\ldots}\inserted{which} admit a dévissage into finite locally free groups, discrete groups, diagonalizable groups, and \deleted{discre} twisted versions\ldots\ If one wants to include the group $\mathbb{G}_a$, one must \ldots

Cartier duality dominates the theory of commutative algebraic (or formal) groups, affine or not -- but it must then be supplemented by consideration of the $\underline{\mathrm{Ext}}^1(G,\mathbb{G}_m)$.
% archive sheet number: 15
\clearpage
